% GENERATED FILE. Edit canonical lessons, then run npm run generate:books.
% canonical-source-hash: fnv1a64:1c55f551e7150c20
% canonical-lessons: SA-C387-unnatih, SA-C387-vikasah, SA-C387-durgatih, SA-C387-samskaroti, SA-C387-vrddhih

\chapter{Getting Better}
\label{ch:sa-getting-better}
\hlchaptermodality{387}

\begin{chapteropening}
\textbf{By the end of this chapter:} \emph{I can say whether things will get better if we work at them.}

Complete the last lesson of chapter 387: I can say whether things will get better if we work at them.
\end{chapteropening}

\section[\texorpdfstring{\sk{उन्नतिः} (unnatiḥ)}{उन्नतिः (unnatiḥ)}]{\sk{उन्नतिः} (unnatiḥ) — advancement; a rise}

\label{lesson:SA-C387-unnatih}



\begin{quote}
Before the new one: say the Sanskrit for a rise in fortune, then the Sanskrit for prosperity.
\end{quote}

\subsection*{You'll want to know: \sk{उन्नतिः}}
\textbf{\sk{उन्नतिः}} — \emph{unnatiḥ} — \textquotedblleft{}advancement; a rise\textquotedblright{}.

\begin{quote}
\sk{यदि} \sk{शिक्षा} \sk{उत्तमा}, \sk{तर्हि} \sk{समाजस्य} \sk{उन्नतिः} \sk{भविष्यति।}
\end{quote}

\begin{quote}
\emph{yadi śikṣā uttamā, tarhi samājasya unnatiḥ bhaviṣyati.}
\end{quote}

\begin{quote}
If education is good, society will advance.
\end{quote}

A feminine noun.

\begin{cousinweb}
\textbf{\sk{उन्नतिः}} (\emph{unnatiḥ}) is built on \emph{ut}, \textquotedblleft{}up\textquotedblright{}, and the root √nam, \textquotedblleft{}to bend\textquotedblright{}; MW gives \emph{un-nati} \textquotedblleft{}rising, ascending\textquotedblright{} and \textquotedblleft{}elevation, height\textquotedblright{}, and \textquotedblleft{}advancement\textquotedblright{} is how modern Sanskrit uses it (MW).
\end{cousinweb}

\begin{grammarlens}[title={What you've built}]
One more word for saying what happens if.
\end{grammarlens}

\subsection*{Your turn}
\begin{itemize}
\raggedright
  \item \emph{Say it:} \emph{unnatiḥ}
  \item \emph{Say it:} \emph{unnatiḥ}, once more
  \item \emph{Say it:} \emph{yadi śikṣā uttamā, tarhi samājasya unnatiḥ bhaviṣyati.} — If education is good, society will advance.
  \item \emph{Recall:} say the Sanskrit for a rise in fortune, then the Sanskrit for prosperity, then say \emph{unnatiḥ} again
\end{itemize}

\begin{tcolorbox}[breakable,colback=teal!4,colframe=teal!35!black,title={Before you move on}]
What does \sk{उन्नतिः} mean? (\textquotedblleft{}Advancement; a rise\textquotedblright{}.) Say it once more.
\end{tcolorbox}
\section[\texorpdfstring{\sk{विकासः} (vikāsaḥ)}{विकासः (vikāsaḥ)}]{\sk{विकासः} (vikāsaḥ) — development; growth}

\label{lesson:SA-C387-vikasah}



\begin{quote}
Before the new one: say the Sanskrit for prosperity, then the Sanskrit for advancement.
\end{quote}

\subsection*{You'll want to know: \sk{विकासः}}
\textbf{\sk{विकासः}} — \emph{vikāsaḥ} — \textquotedblleft{}development; growth\textquotedblright{}.

\begin{quote}
\sk{यदि} \sk{मार्गाः} \sk{भवन्ति}, \sk{तर्हि} \sk{ग्रामस्य} \sk{विकासः} \sk{भविष्यति।}
\end{quote}

\begin{quote}
\emph{yadi mārgāḥ bhavanti, tarhi grāmasya vikāsaḥ bhaviṣyati.}
\end{quote}

\begin{quote}
If there are roads, the village will develop.
\end{quote}

A masculine noun.

\begin{cousinweb}
\textbf{\sk{विकासः}} (\emph{vikāsaḥ}) is built on \emph{vi} and the root √kas, \textquotedblleft{}to open\textquotedblright{}; MW gives \emph{vi-kāsa} \textquotedblleft{}expanding, budding\textquotedblright{} and \textquotedblleft{}expansion, development, growth\textquotedblright{} (MW).
\end{cousinweb}

\begin{grammarlens}[title={What you've built}]
One more word for saying what happens if.
\end{grammarlens}

\subsection*{Your turn}
\begin{itemize}
\raggedright
  \item \emph{Say it:} \emph{vikāsaḥ}
  \item \emph{Say it:} \emph{vikāsaḥ}, once more
  \item \emph{Say it:} \emph{yadi mārgāḥ bhavanti, tarhi grāmasya vikāsaḥ bhaviṣyati.} — If there are roads, the village will develop.
  \item \emph{Recall:} say the Sanskrit for prosperity, then the Sanskrit for advancement, then say \emph{vikāsaḥ} again
\end{itemize}

\begin{tcolorbox}[breakable,colback=teal!4,colframe=teal!35!black,title={Before you move on}]
What does \sk{विकासः} mean? (\textquotedblleft{}Development; growth\textquotedblright{}.) Say it once more.
\end{tcolorbox}
\section[\texorpdfstring{\sk{दुर्गतिः} (durgatiḥ)}{दुर्गतिः (durgatiḥ)}]{\sk{दुर्गतिः} (durgatiḥ) — hard times; misfortune}

\label{lesson:SA-C387-durgatih}



\begin{quote}
Before the new one: say the Sanskrit for advancement, then the Sanskrit for development.
\end{quote}

\subsection*{You'll want to know: \sk{दुर्गतिः}}
\textbf{\sk{दुर्गतिः}} — \emph{durgatiḥ} — \textquotedblleft{}hard times; misfortune\textquotedblright{}.

\begin{quote}
\sk{यदि} \sk{कार्यं} \sk{न} \sk{अस्ति}, \sk{तर्हि} \sk{ग्रामस्य} \sk{दुर्गतिः} \sk{भवति।}
\end{quote}

\begin{quote}
\emph{yadi kāryaṃ na asti, tarhi grāmasya durgatiḥ bhavati.}
\end{quote}

\begin{quote}
If there is no work, the village falls on hard times.
\end{quote}

A feminine noun.

\begin{cousinweb}
\textbf{\sk{दुर्गतिः}} (\emph{durgatiḥ}) is built from \emph{dur-}, \textquotedblleft{}bad, hard\textquotedblright{}, and \emph{gati}, \textquotedblleft{}a going, a way\textquotedblright{}, from the root √gam, \textquotedblleft{}to go\textquotedblright{}; MW gives \emph{dur-gati} \textquotedblleft{}misfortune, distress, poverty\textquotedblright{} (MW).
\end{cousinweb}

\begin{grammarlens}[title={What you've built}]
One more word for saying what happens if.
\end{grammarlens}

\subsection*{Your turn}
\begin{itemize}
\raggedright
  \item \emph{Say it:} \emph{durgatiḥ}
  \item \emph{Say it:} \emph{durgatiḥ}, once more
  \item \emph{Say it:} \emph{yadi kāryaṃ na asti, tarhi grāmasya durgatiḥ bhavati.} — If there is no work, the village falls on hard times.
  \item \emph{Recall:} say the Sanskrit for advancement, then the Sanskrit for development, then say \emph{durgatiḥ} again
\end{itemize}

\begin{tcolorbox}[breakable,colback=teal!4,colframe=teal!35!black,title={Before you move on}]
What does \sk{दुर्गतिः} mean? (\textquotedblleft{}Hard times; misfortune\textquotedblright{}.) Say it once more.
\end{tcolorbox}
\section[\texorpdfstring{\sk{संस्करोति} (saṃskaroti)}{संस्करोति (saṃskaroti)}]{\sk{संस्करोति} (saṃskaroti) — refines; improves}

\label{lesson:SA-C387-samskaroti}



\begin{quote}
Before the new one: say the Sanskrit for development, then the Sanskrit for hard times.
\end{quote}

\subsection*{You'll want to know: \sk{संस्करोति}}
\textbf{\sk{संस्करोति}} — \emph{saṃskaroti} — \textquotedblleft{}refines; improves\textquotedblright{}.

\begin{quote}
\sk{यदि} \sk{लेखकः} \sk{लेखं} \sk{संस्करोति}, \sk{तर्हि} \sk{सः} \sk{उत्तमः} \sk{भविष्यति।}
\end{quote}

\begin{quote}
\emph{yadi lekhakaḥ lekhaṃ saṃskaroti, tarhi saḥ uttamaḥ bhaviṣyati.}
\end{quote}

\begin{quote}
If the writer polishes his article, it will be excellent.
\end{quote}

Past (he …): \sk{सः} \sk{लेखं} \sk{संस्कृतवान्}, \emph{saḥ lekhaṃ saṃskṛtavān}.

\begin{cousinweb}
\textbf{\sk{संस्करोति}} (\emph{saṃskaroti}) is built on \emph{sam}, \textquotedblleft{}together\textquotedblright{}, and the root √kṛ, \textquotedblleft{}to make\textquotedblright{}, with an inserted \emph{s}; MW gives \emph{saṃ-s-√kṛ} \textquotedblleft{}to put together, to prepare, to refine\textquotedblright{} and \emph{saṃ-skṛta} \textquotedblleft{}perfected\textquotedblright{} (MW).
\end{cousinweb}

\begin{grammarlens}[title={What you've built}]
One more word for saying what happens if.
\end{grammarlens}

\subsection*{Your turn}
\begin{itemize}
\raggedright
  \item \emph{Say it:} \emph{saṃskaroti}
  \item \emph{Say it:} \emph{saṃskaroti}, once more
  \item \emph{Say it:} \emph{yadi lekhakaḥ lekhaṃ saṃskaroti, tarhi saḥ uttamaḥ bhaviṣyati.} — If the writer polishes his article, it will be excellent.
  \item \emph{Recall:} say the Sanskrit for development, then the Sanskrit for hard times, then say \emph{saṃskaroti} again
\end{itemize}

\begin{tcolorbox}[breakable,colback=teal!4,colframe=teal!35!black,title={Before you move on}]
What does \sk{संस्करोति} mean? (\textquotedblleft{}Refines; improves\textquotedblright{}.) Say it once more.
\end{tcolorbox}
\section[\texorpdfstring{\sk{वृद्धिः} (vṛddhiḥ)}{वृद्धिः (vṛddhiḥ)}]{\sk{वृद्धिः} (vṛddhiḥ) — growth; increase}

\label{lesson:SA-C387-vrddhih}



\begin{quote}
Before the new one: say the Sanskrit for hard times, then the Sanskrit for refines.
\end{quote}

\subsection*{You'll want to know: \sk{वृद्धिः}}
\textbf{\sk{वृद्धिः}} — \emph{vṛddhiḥ} — \textquotedblleft{}growth; increase\textquotedblright{}.

\begin{quote}
\sk{यदि} \sk{वर्षा} \sk{उत्तमा}, \sk{तर्हि} \sk{धान्यस्य} \sk{वृद्धिः} \sk{भविष्यति।}
\end{quote}

\begin{quote}
\emph{yadi varṣā uttamā, tarhi dhānyasya vṛddhiḥ bhaviṣyati.}
\end{quote}

\begin{quote}
If the rain is good, the grain will grow well.
\end{quote}

A feminine noun.

\begin{cousinweb}
\textbf{\sk{वृद्धिः}} (\emph{vṛddhiḥ}) comes from the root √vṛdh, \textquotedblleft{}to grow\textquotedblright{}; MW gives \emph{vṛddhi} \textquotedblleft{}growth, increase, augmentation, rise\textquotedblright{} (MW).
\end{cousinweb}

\begin{grammarlens}[title={What you've built}]
One more word for saying what happens if.
\end{grammarlens}

\subsection*{Your turn}
\begin{itemize}
\raggedright
  \item \emph{Say it:} \emph{vṛddhiḥ}
  \item \emph{Say it:} \emph{vṛddhiḥ}, once more
  \item \emph{Say it:} \emph{yadi varṣā uttamā, tarhi dhānyasya vṛddhiḥ bhaviṣyati.} — If the rain is good, the grain will grow well.
  \item \emph{Recall:} say the Sanskrit for hard times, then the Sanskrit for refines, then say \emph{vṛddhiḥ} again
\end{itemize}

\begin{tcolorbox}[breakable,colback=teal!4,colframe=teal!35!black,title={Before you move on}]
What does \sk{वृद्धिः} mean? (\textquotedblleft{}Growth; increase\textquotedblright{}.) Say it once more.
\end{tcolorbox}
